% Generated draft: requires mathematical and source review.
\documentclass[11pt]{article}
\usepackage[margin=25mm]{geometry}
\usepackage{fontspec}
\setmainfont{Latin Modern Roman}
\usepackage{amsmath,amssymb,mathtools,unicode-math}
\setmathfont{Latin Modern Math}
\usepackage{enumitem,needspace,xurl,hyperref}
\setlength{\parindent}{0pt}
\setlength{\parskip}{5pt}
\setlength{\emergencystretch}{3em}
\newcommand{\N}{\mathbb{N}}
\newcommand{\Z}{\mathbb{Z}}
\newcommand{\Q}{\mathbb{Q}}
\newcommand{\R}{\mathbb{R}}
\newcommand{\C}{\mathbb{C}}
\newcommand{\F}{\mathbb{F}}
\newcommand{\A}{\mathbb{A}}
\newcommand{\PP}{\mathbb{P}}
\newcommand{\E}{\mathbb{E}}
\newcommand{\eps}{\varepsilon}
\newcommand{\dd}{\,\mathrm{d}}
\begin{document}
\section*{UnsolvedMath ID: 701490}
\textbf{Source problem code:} catalogue entry
\subsection{Definitions and mathematical statement}
A \emph{\(2\)-group} is understood here in the usual periodic sense: a group \(G\), possibly infinite, such that every element has finite order equal to a power of \(2\). No local finiteness assumption is imposed. For a group \(H\), define its lower central series by
\[
\gamma_1(H)=H,\qquad \gamma_{n+1}(H)=[\gamma_n(H),H]\quad(n\geq 1),
\]
where \([A,B]\) denotes the subgroup generated by the commutators \(a^{-1}b^{-1}ab\) with \(a\in A\) and \(b\in B\). Thus \(H\) is nilpotent of class at most \(c\) if \(\gamma_{c+1}(H)=\{1\}\).

\textbf{Question (D. V. Lytkina, Kourovka Notebook 18.61).} Is it true that every \(2\)-group \(G\) satisfying
\[
\forall H\leq G,\qquad |H|<\infty\ \Longrightarrow\ \gamma_4(H)=\{1\}
\]
is nilpotent; equivalently, must there exist an integer \(c\geq 1\) such that \(\gamma_{c+1}(G)=\{1\}\)? The integer \(c\) may depend on \(G\); the question does not require \(G\) itself to have class at most \(3\).

The supplied record additionally states that the answer is affirmative when the hypothesis on finite subgroups is replaced by \(\gamma_3(H)=\{1\}\), meaning class at most \(2\). This is retained as a supplied background assertion; no separate supporting citation is supplied. The record does not explicitly define its convention for an infinite \(2\)-group.
\subsection{Short English statement}
Must a possibly infinite \(2\)-group be nilpotent if every finite subgroup is nilpotent of class at most \(3\)? The source states that the analogous implication for class at most \(2\) holds. Its label of open is supplied registry information, not an independently verified assessment.
\subsection{Sources}
\begin{itemize}
\item Supplied local UnsolvedMath registry record, permanent ID \(701490\), exact source code \texttt{catalogue entry}, title \emph{Kourovka Notebook Problem 18.61}; upstream revision \texttt{d371ee917ac9647fbb06cd6fc529deec6f3dda82}. The record was created and updated on 2026-10-16 and labels the problem open.
\item \emph{The Kourovka Notebook (Unsolved Problems in Group Theory)}, problem 18.61, proposed by D. V. Lytkina in 2014. The supplied background identifies the 18th issue (2014) and the 21st issue edition (2026), arXiv:1401.0300v48. The collection description names E. I. Khukhro and V. D. Mazurov as editors of Issues 1--21 (1965--2026). These bibliographic details are supplied attribution and have not been independently checked.
\end{itemize}
\end{document}
